%% ma: L11-B06-0025
%% lop: 11
%% chuong: 2
%% bai: 6
%% dang: 11.6.10
%% loai: TL
%% muc: VD
%% dap_an: "-"
%% nguon: "de-on-tap-giua-ki-1-toan-11-nam-2026-2027-truong-thpt-ngo-quyen-da-nang.pdf (M09, Đề ôn tập giữa kì 1 Toán 11 THPT Ngô Quyền, Đà Nẵng), Mã đề 0113, Phần IV Câu 2"
%% kien_thuc: "Bài 6 (cấp số cộng), diện tích và đường cao tam giác đều (lớp 9, lớp 10)"
%% trang_thai: cho-duyet
%% ly_do: "Dạng mới đề xuất 11.6.10"
%% ghi_chu: "Hiểu số hàng là số đường song song với đáy, cách đỉnh 0, 1, 2, ... mét còn nằm trong tam giác."
\begin{ex}
Ông An trồng cây cao su trên một khu đất hình tam giác đều có diện tích $2019\text{ m}^2$. Hàng thứ nhất ông trồng $1$ cây ở đỉnh, hàng thứ hai cách đỉnh $1$ m ông trồng nhiều hơn hàng thứ nhất $2$ cây và cứ tiếp tục như vậy đến hàng cuối cùng (mỗi hàng cách đều nhau $1$ m và song song với cạnh đáy của khu đất). Hỏi ông An trồng được tất cả bao nhiêu cây cao su?
\loigiai{Tam giác đều cạnh $a$ có diện tích $\dfrac{a^2\sqrt3}4$ và đường cao $h=\dfrac{a\sqrt3}2$, nên diện tích bằng $\dfrac{h^2}{\sqrt3}$. Suy ra $h=\sqrt{2019\sqrt3}\approx59{,}14$ (m).

Hàng thứ $k$ cách đỉnh $k-1$ (m), cần $k-1\le59{,}14$ nên có $60$ hàng. Số cây các hàng lập thành cấp số cộng $u_1=1$, $d=2$: $S_{60}=\dfrac{60}2(2+59\cdot2)=3600$ (cây).}
\end{ex}
